% Gauss, Werke X/1 (1917), native leaf44, printed37.
\SourceLeaf{044}\NoSourceHead
\HistoricTitle{[KUBISCHE UND BIQUADRATISCHE RESTE.]}
\SourceGroup{V}
\HistoricTitle{[ZWEI SÄTZE ÜBER KUBISCHE UND BIQUADRATISCHE RESTE.]}
\HistoricCitation{[Aus Handbuch 18, Bd. Mathematische Brouillons, October 1805, Einbanddecke.]}
\par\indent Wenn eine Primzahl $p$ sich unter die Form
\[
aa+27bb
\]
\noindent setzen lässt, so ist $2$ Res. Cub. ipsius $p$.

Wenn sich dieselbe unter die Form
\[
aa+64bb
\]
\noindent setzen lässt, so ist $\pm2$ ein Biquadr. Rest von $p$.
\begin{center}\textbf{BEMERKUNG.}\end{center}
\begingroup\small
\par\indent Vergl. hierzu die Werke VIII, S. 6 abgedruckte Aufzeichnung sowie die \textit{Theoria Residuorum Biquadraticorum, Commentatio secunda} (1832) art. 24, Werke II, S. 95, 96.
\begin{flushright}\textsc{Schlesinger.}\end{flushright}
\ClosingRule\clearpage\endgroup
